% ============================================================
% Section 3: Methodology
% ============================================================

\section{Methodology}
\label{sec:methodology}

\subsection{Problem Formulation}
\label{subsec:formulation}

We consider the dimensionless square domain $\Omega = [0,1]^2$, where the
coordinates have been normalized by the laser wavelength
($\lambda = 638\,\mathrm{nm}$).
Under this normalization, the dimensionless wavenumber is $\tilde{k} = 2\pi$.

\paragraph{1D validation.}
The one-dimensional Helmholtz equation on $x \in [0,1]$ with Dirichlet conditions:
%
\begin{equation}
\label{eq:helmholtz1d}
\begin{aligned}
  &\frac{d^2 E}{dx^2} + k^2 E = 0, \\
  &E(x_j) = \cos(k x_j), \\
  &x_j \in \{0,\ 0.25,\ 0.5,\ 0.75,\ 1\}
\end{aligned}
\end{equation}
%
admits the analytical solution $E_{\mathrm{exact}}(x) = \cos(kx)$, used to validate
the SIREN architecture. Five boundary points are imposed instead of only the two
endpoints: the symmetric condition $E(0) = E(1) = 1$ collapses the network towards
the trivial constant solution during training, since both endpoints are consistent
with $E \equiv 1$; the three intermediate points break that symmetry and stabilize
convergence.

\paragraph{2D validation.}
The 2D Helmholtz equation on $\Omega = [0,1]^2$ with a plane-wave condition on the
four edges:
%
\begin{equation}
  E_{\mathrm{exact}}(x,y) = e^{i(k_x x + k_y y)}, \quad k_x = k_y = \frac{k}{\sqrt{2}},
  \label{eq:solution2d}
\end{equation}
%
which satisfies $k_x^2 + k_y^2 = k^2$.

\subsection{SIREN Architecture}
\label{subsec:architecture}

We employ a SIREN architecture with parameters consolidated from the 1D validation:
$L = 5$ hidden layers, dimension $d = 128$ (doubled with respect to the 1D case in
order to accommodate the complex field), input frequency $\omega_0 = 1.0$ and a
linear output layer with two neurons
$(E_{\mathrm{real}}, E_{\mathrm{imag}})$.
The model has 66{,}690 trainable parameters.

The choice $\omega_0 = 1.0$, instead of the value $\omega_0 = 30$ recommended by
Sitzmann et al. \cite{Sitzmann2020_SIREN} for image signals, is justified by the
nature of the problem: the dimensionless domain $[0,1]^2$ with $k = 2\pi$ does not
require capturing high frequencies, and larger values of $\omega_0$ destabilize
training with L-BFGS.

\subsection{Latin Hypercube Sampling}
\label{subsec:lhs}

In the 2D case, the $N_c = 3{,}000$ collocation points are generated through Latin
Hypercube Sampling (LHS) \cite{Mckay1979_LHS} instead of a uniform grid.
LHS divides each dimension into $N_c$ equal strata and places exactly one point per
stratum, guaranteeing uniform coverage of the domain without the redundant patterns
of a Cartesian grid.
Additionally, the $N_b = 300$ boundary points per edge are uniformly distributed
along each side of the domain.

\subsection{Loss Function and Weights}
\label{subsec:loss}

The total loss function for the 2D case is:
%
\begin{equation}
  \mathcal{L} = \mathcal{L}_{\mathrm{data}}
              + \lambda_{\mathrm{phys}}\!\left(
                  \mathcal{L}_{R_{\mathrm{real}}} + \mathcal{L}_{R_{\mathrm{imag}}}
                \right),
  \label{eq:loss_2d}
\end{equation}
%
with physics weight $\lambda_{\mathrm{phys}} = 0.1$.
This empirically calibrated value balances the fit to the boundary data and the
fulfillment of the PDE: larger values ($\lambda = 1.0$) overweight the physics and
prevent convergence in 2D; smaller values produce fields that violate the Helmholtz
equation.
Table~\ref{tab:ablation_lambda} quantifies the effect on the average $L^2$ error.
This fragility with respect to the relative weight of the loss terms is not an
isolated phenomenon of this implementation: Krishnapriyan et al.
\cite{Krishnapriyan2021_failureModes} systematically document that PINNs with fixed
weights between the data term and the physical residual term may fail to converge
even when the global loss appears to decrease, attributing the problem to an
optimization landscape ill-conditioned by the competition between both terms.
\begin{table}[t]
\centering
\footnotesize
\renewcommand{\arraystretch}{1.2}
\setlength{\tabcolsep}{4pt}
\caption{ablation of the physics weight $\lambda_{\mathrm{phys}}$, NB02}
\label{tab:ablation_lambda}
\begin{tabular}{@{}c c c l@{}}
\hline
$\lambda_{\mathrm{phys}}$ & Avg.\ $L^2$ error (\%) & Converges & Remark \\
\hline
0.01  & 0.163 & Yes & Physics underweighted \\
0.1   & 0.171 & Yes & \textbf{Optimal config.} \\
1.0   & $>5.0$ & No & Physics overweighted \\
\hline
\end{tabular}
\end{table}
\noindent\footnotesize SIREN $5\times128$ architecture, $\omega_0 = 1.0$, SEED = 42.
Same Adam + L-BFGS protocol as NB02.
``Does not converge'' indicates that the Adam optimizer did not reach the stopping
condition before 15{,}000 epochs with $\varepsilon_{L^2} < 5\%$.
\normalsize

\subsection{Optimization Strategy}
\label{subsec:optimization}

Training follows a two-phase strategy:

\paragraph{Phase 1 - Adam.}
Adam optimizer with learning rate $\eta = 10^{-3}$, up to 15{,}000 epochs, with
early stopping after a patience of 800 epochs without improvement greater than
$\delta = 10^{-7}$.
Adam performs a global exploration of the parameter space and provides a starting
point close to the global minimum.

\paragraph{Phase 2 - L-BFGS.}
L-BFGS optimizer with \textit{strong Wolfe} line search, a maximum of 1{,}000
iterations and history size 100.
L-BFGS exploits the curvature information of the loss landscape to perform a
high-precision refinement from the point delivered by Adam.

This combination, a widespread practice in the PINN literature
\cite{Cuomo2022_review}, outperforms pure Adam in final accuracy and pure L-BFGS in
robustness against local minima.

\subsection{Validation Metric}
\label{subsec:metrics}

\paragraph{Relative L2 error} (NB01 and NB02):
%
\begin{equation}
  \varepsilon_{L2} = \frac{\|\hat{E} - E_{\mathrm{exact}}\|_2}{\|E_{\mathrm{exact}}\|_2},
  \label{eq:l2}
\end{equation}
%
evaluated on a uniform $100 \times 100$ grid.
There is no standardized error tolerance criterion in the PINN literature for this
type of validation, so $\varepsilon_{L2} < 5\%$ was adopted as an engineering
threshold for this work. This value is, moreover, considerably lower than the order
of magnitude reported in comparable PINN implementations for the Helmholtz equation
(Table~\ref{tab:comparison}).

\subsection{Implementation}
\label{subsec:implementation}

The complete implementation is written in Python with PyTorch and executed on an
NVIDIA RTX 5050 GPU (8.5 GB VRAM).
The modules \texttt{src/models.py}, \texttt{src/losses.py} and
\texttt{src/utils.py} contain the reusable classes and functions; the notebooks
reproduce the experiments in a self-contained manner.
The global seed \texttt{SEED} = 42 guarantees reproducibility.
